% Apuntes de teoría en PDF · 4.º ESO · Unidad 1: Naturales, enteros y fracciones (versión en castellano)
% Mateix estil que els apunts de Batxillerat (franja i caixes de teoria en blau, exemples en roig), però més visuals,
% curts i amb dibuixos, per a l'alumnat de 3r ESO. Es compila amb pdflatex dues vegades (per als enllaços de l'índex),
% amb logo-fentmates.pdf (_eines/logo/pdf/horitzontal-clar.pdf) al costat.
\documentclass[11pt,a4paper]{article}
\usepackage[utf8]{inputenc}
\usepackage[T1]{fontenc}
\usepackage[spanish,es-noshorthands,es-nodecimaldot,es-nolists]{babel}
\usepackage[sfdefault,book,lining]{FiraSans}
\usepackage[italic,noplusnominus,noequal]{mathastext}
\linespread{1.18}
\usepackage[a4paper,top=2.4cm,bottom=2.2cm,left=1.6cm,right=1.6cm,headheight=1.1cm,headsep=0.5cm,footskip=1cm]{geometry}
\usepackage{amsmath,amssymb,cancel}
\usepackage{icomma}
\usepackage{graphicx}
\usepackage[table]{xcolor}
\usepackage{tikz}\usetikzlibrary{calc,positioning,arrows.meta,decorations.pathreplacing,shapes.geometric}
\usepackage{fontawesome5}
\usepackage{ccicons}
\usepackage{fancyhdr}
\usepackage{enumitem}
\usepackage{array}
\usepackage{needspace}
\usepackage[most]{tcolorbox}
\usepackage{titlesec}
\usepackage[hidelinks]{hyperref}
\hypersetup{pdftitle={Apuntes de teoría. Unidad 1: Naturales, enteros y fracciones. 4.º ESO},pdfsubject={Matemáticas A 4.º ESO},pdfkeywords={números naturales, división entera, potencias, raíz cuadrada, jerarquía de las operaciones, números enteros, valor absoluto, fracciones, operaciones combinadas, problemas con fracciones, propiedades de las potencias, exponente negativo}}

% colors de la web (bloc d'Àlgebra: blau)
\definecolor{fmnavy}{HTML}{1B2748}
\definecolor{fmblue}{HTML}{3D6AA0}
\definecolor{fmbluel}{HTML}{4E81BD}
\definecolor{fmyellow}{HTML}{FFC94D}
\definecolor{fmgold}{HTML}{E9A919}
\definecolor{fmaccent}{HTML}{FF7A50}
\definecolor{fmdef}{HTML}{1D4F73}\definecolor{fmdefbg}{HTML}{EEF4F9}
\definecolor{fmex}{HTML}{A3203A}\definecolor{fmexbg}{HTML}{FBF0F2}
\definecolor{fmsf}{HTML}{C2410C}
\definecolor{fmok}{HTML}{15803D}
\definecolor{fmmuted}{HTML}{5D6685}
\definecolor{fmsurface}{HTML}{EEF1F8}
\definecolor{fmA}{HTML}{2F6FD6}\definecolor{fmB}{HTML}{EA6A10}\definecolor{fmR}{HTML}{E11D48}\definecolor{fmK}{HTML}{16A34A}\definecolor{fmV}{HTML}{7C3AED}

% lletres
\newcommand{\comic}{\fontfamily{ComicNeue-TLF}\selectfont}
\newcommand{\heavy}{\fontseries{eb}\selectfont}
\newcommand{\semi}{\fontseries{sb}\selectfont}
\DeclareUnicodeCharacter{2212}{-}

\setlength{\parindent}{0pt}
\setlength{\parskip}{7pt}
\setlength{\lineskip}{4pt}\setlength{\lineskiplimit}{4pt}
\setlength{\emergencystretch}{3em}
\setlist{itemsep=3pt,topsep=3pt,leftmargin=1.4em}
\setcounter{secnumdepth}{1}
\renewcommand{\thesection}{1.\arabic{section}}
\titleformat{\section}[hang]{\color{fmnavy}\fontsize{23}{27}\heavy\selectfont}{\tcbox[on line,colback=fmblue,colframe=fmblue,arc=3mm,boxsep=2pt,left=6pt,right=6pt,top=3pt,bottom=3pt]{\color{white}\thesection}}{0.45em}{}[\vspace{5pt}{\color{fmblue}\titlerule[2pt]}]
\titlespacing*{\section}{0pt}{4pt}{6pt}
\titleformat{\subsection}{\color{fmblue}\Large\semi}{}{0pt}{\faCaretRight\ }
\titlespacing*{\subsection}{0pt}{8pt}{2pt}

% frase de cada apartat, «a mà»
\newcommand{\frase}[1]{\par\vspace{-2pt}{\comic\Large\color{fmbluel}\faCommentDots[regular]\ \ #1\par}\vspace{2pt}}
\newcommand{\seccio}[3]{\Needspace{16\baselineskip}\section{#2}\hypertarget{#1}{}\frase{#3}}
% notes a mà dins dels exemples
\newcommand{\nota}[1]{{\comic\small\color{fmex!85!black}#1}}
\newcommand{\notab}[1]{{\comic\small\color{fmblue}#1}}
% bé / malament
\newcommand{\be}[1]{{\color{fmok}\faCheckCircle}\ \ensuremath{#1}}
\newcommand{\mal}[1]{{\color{fmR}\faTimesCircle}\ {\color{fmR}\ensuremath{#1}}}
% pas numerat
\newcommand{\pas}[1]{\tikz[baseline=(p.base)]\node[circle,fill=fmblue,text=white,inner sep=1.2pt,minimum size=1.25em,font=\small\heavy](p){#1};}
% etiqueta de color
\newcommand{\chip}[2]{\tcbox[on line,colback=#1!14,colframe=#1!60,arc=2mm,boxsep=1pt,left=5pt,right=5pt,top=1pt,bottom=1pt]{#2}}
\newcommand{\fmkey}[1]{\colorbox{fmyellow!55}{\strut #1}}
\newcommand{\cA}[1]{{\color{fmA}#1}}\newcommand{\cB}[1]{{\color{fmB}#1}}\newcommand{\cK}[1]{{\color{fmK}#1}}\newcommand{\cR}[1]{{\color{fmR}#1}}\newcommand{\cV}[1]{{\color{fmV}#1}}

% caixes
\tcbset{fmbase/.style={enhanced,arc=3mm,boxrule=1.2pt,left=4mm,right=4mm,top=2mm,bottom=2.5mm,before skip=8pt,after skip=8pt,fonttitle=\heavy\large,before upper=\raggedright}}
\newtcolorbox{idea}[1][Idea clave]{fmbase,colback=fmdefbg,colframe=fmdef,coltitle=white,title={\faKey\ \ #1}}
\newtcolorbox{exemple}[1][Ejemplo]{fmbase,colback=fmexbg,colframe=fmex,colbacktitle=fmexbg,coltitle=fmex,title={\faPencil*\ \ #1},titlerule=0.5pt,titlerule style={fmex!40,dashed}}
\newtcolorbox{compte}[1][¡Cuidado!]{fmbase,colback=fmsf!6,colframe=fmsf,coltitle=white,fonttitle=\comic\bfseries\Large,title={\faExclamationTriangle\ \ #1}}
\newtcolorbox{truc}[1][Truco]{fmbase,colback=fmyellow!20,colframe=fmgold,colbacktitle=fmyellow,coltitle=fmnavy,fonttitle=\comic\bfseries\Large,title={\faLightbulb[regular]\ \ #1}}

% dibuix de la targeta de la unitat, en blanc
\newcommand{\fmcard}{\begin{tikzpicture}[x=0.062cm,y=-0.062cm,white,line width=3pt,line cap=round]
  \begin{scope}[transparency group,opacity=0.15]
  \node[anchor=base west,inner sep=0pt,font=\fontsize{30}{30}\selectfont] at (14,44) {$-3\quad\frac34\quad 2^{-1}$};
  \draw[line width=2.4pt] (20,95) -- (230,95);
  \foreach \x in {40,75,110,145,180,215} \draw[line width=2.4pt] (\x,88) -- (\x,102);
  \end{scope}
\end{tikzpicture}}
\newcommand{\fmbrand}{\raisebox{-0.25\height}{\includegraphics[height=0.5cm]{logo-fentmates.pdf}}}

% regla de Ruffini: #1 = a, #2 = coeficients, #3 = fila del mig (sense el primer), #4 = fila de baix, #5 = nombre de coeficients
\newcommand{\ruffini}[5]{\begin{tikzpicture}[x=1.25cm,y=0.85cm,every node/.style={font=\large}]
  \foreach \v [count=\i] in {#2} \node (t\i) at (\i,0) {$\v$};
  \foreach \v [count=\i] in {#3} {\pgfmathtruncatemacro\j{\i+1}\node (m\j) at (\j,-1) {$\v$};}
  \foreach \v [count=\i] in {#4} \node (b\i) at (\i,-2.1) {$\v$};
  \node[text=fmB] at (0,-1) {$#1$};
  \draw[line width=1.1pt,fmnavy] (0.45,0.5)--(0.45,-1.5) (0.45,-1.5)--(#5+0.5,-1.5);
  \node[draw=fmR,line width=1.3pt,rounded corners=2mm,minimum size=0.75cm] at (#5,-2.1) {};
  \draw[-{Latex[length=2.2mm]},fmK,line width=1.1pt] (t1.south) -- (b1.north);
  \pgfmathtruncatemacro\last{#5-1}
  \foreach \i in {1,...,\last} {\pgfmathtruncatemacro\j{\i+1}
    \draw[-{Latex[length=2.2mm]},fmB,line width=1.1pt,shorten >=1pt] (b\i.north east) -- (m\j.south west);
    \draw[-{Latex[length=2.2mm]},fmA!70,line width=1pt,dashed,shorten >=2pt] ([xshift=8pt]t\j.south) -- ([xshift=8pt]b\j.north);}
\end{tikzpicture}}

% capçalera i peu
\pagestyle{fancy}\fancyhf{}
\fancyhead[L]{\small\fmbrand}
\fancyhead[R]{\small\color{fmmuted}4.º ESO · Unidad 1: Naturales, enteros y fracciones}
\fancyfoot[L]{\scriptsize\color{fmmuted}\href{https://creativecommons.org/licenses/by-nc-sa/4.0/deed.es}{\ccbyncsa}\ \ FentMates · fentmates.org · CC BY-NC-SA 4.0}
\fancyfoot[R]{\small\color{fmmuted}\thepage}
\renewcommand{\headrulewidth}{0.4pt}\renewcommand{\headrule}{\color{fmmuted!40}\hrule height \headrulewidth}
\fancypagestyle{primera}{\fancyhf{}\renewcommand{\headrulewidth}{0pt}
  \fancyfoot[L]{\scriptsize\color{fmmuted}\href{https://creativecommons.org/licenses/by-nc-sa/4.0/deed.es}{\ccbyncsa}\ \ FentMates · fentmates.org · CC BY-NC-SA 4.0}
  \fancyfoot[R]{\small\color{fmmuted}\thepage}}

% targeta de l'índex
\newcommand{\itile}[4]{\begin{tcolorbox}[enhanced,colback=white,colframe=fmblue!35,arc=3mm,boxrule=1pt,boxsep=0pt,left=6pt,right=6pt,top=5pt,bottom=5pt,before skip=0pt,after skip=7pt]
  \hyperlink{#1}{\begin{minipage}[c]{0.9cm}\centering{\color{fmblue}\LARGE #3}\end{minipage}\hspace{6pt}%
  \begin{minipage}[c]{0.78\linewidth}{\color{fmblue}\heavy #2}\\[-2pt]{\color{fmnavy}\semi\large #4}\end{minipage}}\end{tcolorbox}}

% taula amb vores arrodonides: el contingut es retalla amb la forma (sense @{} a les vores ni entre columnes)
\newsavebox\tabclipbox
\newcommand{\taulaclip}[3][3mm]{\sbox\tabclipbox{#3}\begin{tikzpicture}[baseline=(tcn.base)]
  \begin{scope}\clip[rounded corners=#1] (0,-\dp\tabclipbox) rectangle (\wd\tabclipbox,\ht\tabclipbox);
  \node[inner sep=0pt,outer sep=0pt,anchor=base west] (tcn) at (0,0) {\usebox\tabclipbox};\end{scope}
  \draw[#2,rounded corners=#1] (0,-\dp\tabclipbox) rectangle (\wd\tabclipbox,\ht\tabclipbox);
\end{tikzpicture}}
\newcommand{\taularod}[3][fmblue]{\taulaclip{#1,line width=1.2pt}{\renewcommand{\arraystretch}{1.5}\begin{tabular}{#2}#3\end{tabular}}}
% balança
\newcommand{\balanca}[3]{\begin{tikzpicture}[scale=#3,transform shape,every node/.style={anchor=center}]
  \fill[fmnavy] (-0.45,-1.9) rectangle (0.45,-1.75);
  \fill[fmnavy!80] (-0.12,-1.75) rectangle (0.12,0);
  \fill[fmgold] (0,0.05) circle (0.13);
  \draw[fmnavy,line width=2.4pt,line cap=round] (-2,0) -- (2,0);
  \foreach \s in {-1,1} {
    \draw[fmmuted,line width=0.8pt] (\s*2,0) -- (\s*2-0.75,-1.05) (\s*2,0) -- (\s*2+0.75,-1.05);
    \fill[fmyellow!70,draw=fmgold,line width=1pt] (\s*2-0.95,-1.05) arc[start angle=180,end angle=360,x radius=0.95,y radius=0.32] -- cycle;}
  \node[font=\Large,fill=fmA!15,draw=fmA,rounded corners=1.5mm,inner sep=3pt] at (-2,-0.68) {$#1$};
  \node[font=\Large,fill=fmB!15,draw=fmB,rounded corners=1.5mm,inner sep=3pt] at (2,-0.68) {$#2$};
\end{tikzpicture}}
% eixos amb quadrícula (dins d'un tikzpicture): xmin, xmax, ymin, ymax
\newcommand{\eixos}[4]{\draw[fmmuted!22,line width=0.4pt] (#1,#3) grid (#2,#4);
  \draw[-{Latex[length=2mm]},fmmuted,line width=0.8pt] (#1-0.2,0)--(#2+0.4,0) node[right,font=\small]{$x$};
  \draw[-{Latex[length=2mm]},fmmuted,line width=0.8pt] (0,#3-0.2)--(0,#4+0.4) node[above,font=\small]{$y$};}
% taula de traduccions (diccionari) amb les vores arrodonides
\newcommand{\dic}[2]{#1 & {\large$#2$}\\}
\newcommand{\dicratio}{0.62}
\newcommand{\dichead}{\begin{tabular}{@{}>{\raggedright}p{\dimexpr\dicratio\linewidth-12pt}>{\centering\arraybackslash}p{\dimexpr\linewidth-\dicratio\linewidth-12pt}@{}}\faComment[regular]\ Lo decimos así\dots & \dots\ así\end{tabular}}
\newenvironment{dicbox}{\begin{tcolorbox}[enhanced,arc=3mm,boxrule=1.2pt,colframe=fmblue,colback=white,coltitle=white,fonttitle=\heavy,
  title={\dichead},left=0pt,right=0pt,top=0pt,bottom=1.5mm,boxsep=0pt,lefttitle=6pt,righttitle=6pt,before skip=0pt,after skip=0pt]\rowcolors{1}{white}{fmsurface}%
  \begin{tabular}{@{\hspace{6pt}}>{\raggedright}p{\dimexpr\dicratio\linewidth-12pt}>{\centering\arraybackslash}p{\dimexpr\linewidth-\dicratio\linewidth-12pt}@{\hspace{6pt}}}}
  {\end{tabular}\end{tcolorbox}}
% resum
\newcommand{\rtile}[3]{\begin{tcolorbox}[enhanced,colback=white,colframe=fmblue!45,arc=3mm,boxrule=1pt,left=3mm,right=3mm,top=2.5mm,bottom=2.5mm,fontupper=\large,
  title={\tcbox[on line,colback=fmblue,colframe=fmblue,arc=2mm,boxsep=0pt,left=4pt,right=4pt,top=1pt,bottom=1pt]{\color{white}\small\heavy #1}\ \ #2},coltitle=fmnavy,colbacktitle=fmdefbg,fonttitle=\semi,toptitle=0.5mm,bottomtitle=0.5mm]#3\end{tcolorbox}}
% portada: #1 subtítol a mà
\newcommand{\portada}[1]{\thispagestyle{primera}
\begin{tikzpicture}[remember picture,overlay]
  \fill[fmblue] (current page.north west) rectangle ([yshift=-9.4cm]current page.north east);
  \begin{scope}
    \clip (current page.north west) rectangle ([yshift=-9.4cm]current page.north east);
    \node[anchor=north west,inner sep=0pt] at ([xshift=6.6cm,yshift=-1.3cm]current page.north west) {\fmcard};
  \end{scope}
  \node[anchor=north west,inner sep=0pt] at ([xshift=1.6cm,yshift=-1.1cm]current page.north west) {\tcbox[colback=white,colframe=white,arc=2.5mm,boxsep=0pt,left=7pt,right=7pt,top=5pt,bottom=5pt]{\includegraphics[height=0.75cm]{logo-fentmates.pdf}}};
\end{tikzpicture}
\vspace*{0.55cm}
\tcbox[colback=white!15!fmblue,colframe=white!15!fmblue,arc=2.5mm,boxsep=0pt,left=6pt,right=6pt,top=3pt,bottom=3pt]{\color{white}\small\bfseries Apuntes de teoría · 4.º ESO · Matemáticas A}

\vspace{2pt}
{\raggedright\color{white}\fontsize{30}{36}\heavy\selectfont{\color{fmyellow}Unidad 1}\\[2pt]Naturales, enteros y fracciones\par}
\vspace{8pt}
{\color{white}\comic\Large #1\par}
\vspace*{\dimexpr 7.4cm-\pagetotal\relax}
}
\newcommand{\llegenda}{\par\vspace{16pt}
{\color{fmnavy}\heavy\Large\faPalette\ \ Cómo se leen estos apuntes}\par\smallskip
\begin{tcbraster}[raster columns=4,raster equal height,raster column skip=8pt,raster before skip=12pt]
\tcbox[colback=fmdefbg,colframe=fmdef,arc=3mm,boxrule=1.2pt,halign=center,width=\linewidth]{{\color{fmdef}\Large\faKey}\\[2pt]{\color{fmdef}\heavy Idea clave}\\[-1pt]{\small lo esencial}}
\tcbox[colback=fmexbg,colframe=fmex,arc=3mm,boxrule=1.2pt,halign=center,width=\linewidth]{{\color{fmex}\Large\faPencil*}\\[2pt]{\color{fmex}\heavy Ejemplo}\\[-1pt]{\small paso a paso}}
\tcbox[colback=fmsf!6,colframe=fmsf,arc=3mm,boxrule=1.2pt,halign=center,width=\linewidth]{{\color{fmsf}\Large\faExclamationTriangle}\\[2pt]{\color{fmsf}\heavy ¡Cuidado!}\\[-1pt]{\small errores típicos}}
\tcbox[colback=fmyellow!20,colframe=fmgold,arc=3mm,boxrule=1.2pt,halign=center,width=\linewidth]{{\color{fmgold!80!black}\Large\faLightbulb[regular]}\\[2pt]{\color{fmnavy}\heavy Truco}\\[-1pt]{\small para recordarlo}}
\end{tcbraster}}
\newcommand{\resumtitol}{\clearpage
\hypertarget{resum}{}
{\color{fmnavy}\fontsize{23}{27}\heavy\selectfont\faClipboardList\ \ En resumen}\par\vspace{5pt}{\color{fmblue}\titlerule[2pt]}
\frase{Toda la unidad en una página.}}
% els 5 passos per a resoldre problemes: #2 = text del pas 2
\newcommand{\passosproblema}[2]{\begin{center}
\begin{tikzpicture}[s/.style={draw=fmblue,fill=fmdefbg,rounded corners=3mm,line width=1.2pt,minimum width=2.75cm,minimum height=2.45cm,align=center,font=\small},
  ic/.style={font=\LARGE,text=fmblue}]
  \foreach \i/\ic/\t in {1/\faBookOpen/{\textbf{Leemos}\\y entendemos},2/\faQuestionCircle/{#1},3/\faEquals/{#2},4/\faCogs/{\textbf{Resolvemos}},5/\faCheckDouble/{\textbf{Comprobamos}\\y respondemos}} {
    \node[s] (s\i) at (3.3*\i,0) {};
    \node[ic] at ([yshift=-0.5cm]s\i.north) {\ic};
    \node[align=center,font=\small,anchor=north] at ([yshift=-0.95cm]s\i.north) {\t};
    \node[circle,fill=fmblue,text=white,font=\small\heavy,inner sep=1.5pt,minimum size=1.3em] at (s\i.north west) {\i};}
  \foreach \i [evaluate=\i as \j using int(\i+1)] in {1,...,4} \draw[-{Latex[length=3mm]},fmmuted,line width=1.5pt] (s\i) -- (s\j);
\end{tikzpicture}
\end{center}}

\newcommand{\lo}{\mathopen{]}}\newcommand{\ro}{\mathclose{[}}
% punts ple i buit, nombres dels eixos i taula de valors
\newcommand{\pple}[2][fmblue]{\fill[#1] #2 circle (3.6pt);}
\newcommand{\nums}[4]{\foreach \k in {#1,...,#2} {\ifnum\k=0\else\node[below,font=\scriptsize,text=fmmuted] at (\k,0) {$\k$};\fi}
  \foreach \k in {#3,...,#4} {\ifnum\k=0\else\node[left,font=\scriptsize,text=fmmuted] at (0,\k) {$\k$};\fi}}
\newcommand{\tv}[3][fmblue]{\taulaclip{#1,line width=1.2pt}{\renewcommand{\arraystretch}{1.45}%
  \begin{tabular}{>{\columncolor{#1}\color{white}\heavy}c*{12}{w{c}{0.8cm}}}#2\\\rowcolor{#1!8}\cellcolor{#1}#3\end{tabular}}}
% etiqueta d'una recta o corba
\newcommand{\etq}[3]{\node[fill=white,inner sep=1.5pt,text=#1,font=\small] at #2 {#3};}


\newcommand{\recta}[2]{\draw[-{Latex[length=2.5mm]},fmnavy,line width=1.1pt] (#1-0.6,0) -- (#2+0.6,0);
  \foreach \k in {#1,...,#2} {\draw[fmnavy,line width=0.9pt] (\k,-0.12) -- (\k,0.12); \node[below=3pt,font=\small] at (\k,0) {$\k$};}}
\newcommand{\barra}[4]{\begin{tikzpicture}
  \foreach \i in {1,...,#1} {\pgfmathsetmacro\xa{(\i-1)*#4/#1}\pgfmathsetmacro\xb{\i*#4/#1}
    \ifnum\i>#2 \filldraw[fill=white,draw=#3,line width=1pt] (\xa,0) rectangle (\xb,0.55);
    \else \filldraw[fill=#3!35,draw=#3,line width=1pt] (\xa,0) rectangle (\xb,0.55); \fi}
\end{tikzpicture}}
\def\intvdot#1#2#3{\ifx#2o\filldraw[fill=white,draw=#3,line width=1.3pt] (#1,0) circle (3.3pt);\fi\ifx#2c\fill[#3] (#1,0) circle (3.5pt);\fi}
\newcommand{\intv}[5]{\begin{tikzpicture}[x=0.4cm,baseline=-3pt]
  \draw[fmmuted,line width=0.8pt] (-6,0)--(7,0);
  \foreach \k in {-5,...,6} \draw[fmmuted,line width=0.6pt] (\k,-0.08)--(\k,0.08);
  \ifx#3a\def\xa{-6}\else\def\xa{#1}\fi
  \ifx#4a\def\xb{7}\else\def\xb{#2}\fi
  \ifx#4a\draw[#5,line width=2.6pt,-{Latex[length=2.6mm]}] (\xa,0)--(\xb,0);\else\ifx#3a\draw[#5,line width=2.6pt,{Latex[length=2.6mm]}-] (\xa,0)--(\xb,0);\else\draw[#5,line width=2.6pt] (\xa,0)--(\xb,0);\fi\fi
  \intvdot{#1}{#3}{#5}\intvdot{#2}{#4}{#5}
  \ifx#3a\else\node[below=2pt,font=\scriptsize] at (#1,0) {$#1$};\fi
  \ifx#4a\else\node[below=2pt,font=\scriptsize] at (#2,0) {$#2$};\fi
\end{tikzpicture}}


\begin{document}
\portada{Contar, deber, repartir: los números de cada día y cómo se opera con ellos.}

{\color{fmnavy}\heavy\Large\faMapSigns\ \ ¿Qué aprenderás?}\par\smallskip
\begin{minipage}[t]{0.485\linewidth}
\itile{sec1}{1.1}{\faSortNumericUp}{Los números naturales}
\itile{sec2}{1.2}{\faSuperscript}{Potencias, raíces y jerarquía}
\itile{sec3}{1.3}{\faThermometerHalf}{Enteros. Valor absoluto}
\itile{sec4}{1.4}{\faCalculator}{Operaciones con enteros}
\itile{sec5}{1.5}{\faChartPie}{Fracciones equivalentes}
\itile{sec6}{1.6}{\faPlus}{Operaciones con fracciones}
\end{minipage}\hfill
\begin{minipage}[t]{0.485\linewidth}
\itile{sec7}{1.7}{\faLayerGroup}{Combinadas con fracciones}
\itile{sec8}{1.8}{\faTasks}{Problemas con fracciones}
\itile{sec9}{1.9}{\faIcon{square-root-alt}}{Base entera y racional}
\itile{sec10}{1.10}{\faListOl}{Propiedades de las potencias}
\itile{sec11}{1.11}{\faIcon{minus-circle}}{Exponente cero y negativo}
\itile{resum}{\phantom{1.11}}{\faClipboardList}{En resumen}
\end{minipage}
\llegenda

% ============================================================ 1.1
\clearpage
\seccio{sec1}{Los números naturales}{Para contar y para ordenar.}

\vfill
\begin{minipage}[c]{0.52\linewidth}
\begin{idea}[Los naturales $\mathbb N$]
{\centering\Large$\mathbb N=\{0,\ 1,\ 2,\ 3,\ 4,\ \ldots\}$\par}\smallskip
Sirven para \textbf{contar} (5 cajas, 120 alumnos) y para \textbf{ordenar} (el 3.º clasificado). Empiezan en el 0 y no se acaban nunca: después de $n$ siempre viene $n+1$.
\end{idea}
\end{minipage}\hfill
\begin{minipage}[c]{0.45\linewidth}\centering
\begin{tikzpicture}[x=0.62cm]
  \draw[-{Latex[length=2.4mm]},fmnavy,line width=1.2pt] (0,0) -- (9.6,0);
  \foreach \k in {0,...,9} {\draw[fmnavy,line width=1pt] (\k,-0.12) -- (\k,0.12); \node[below=2pt,font=\small] at (\k,-0.1) {$\k$};}
  \foreach \k in {0,...,8} \draw[fmB,line width=1pt,-{Latex[length=1.6mm]}] (\k+0.1,0.25) to[bend left=50] (\k+0.9,0.25);
  \node[above,font=\scriptsize,text=fmB] at (4.5,0.7) {$+1$ cada paso};
\end{tikzpicture}\\[2pt]
{\comic\small\color{fmmuted}una semirrecta: siempre la misma distancia}
\end{minipage}

\vfill
\begin{minipage}[c]{0.5\linewidth}\centering
\taularod{>{\heavy}l<{\hspace{4pt}}>{\hspace{4pt}}l<{\hspace{4pt}}>{\hspace{4pt}}c<{\hspace{4pt}}}{
\rowcolor{fmblue}\color{white}Operación & \color{white}\heavy Ejemplo & \color{white}\heavy ¿Da un natural?\\
Suma & $7+5=12$ & \be{}\\
\rowcolor{fmsurface}Multiplicación & $7\cdot5=35$ & \be{}\\
Resta & $5-7=\ ?$ & \mal{} \ {\small (enteros)}\\
\rowcolor{fmsurface}División & $7:5=\ ?$ & \mal{} \ {\small (si no es exacta)}
}
\end{minipage}\hfill
\begin{minipage}[c]{0.47\linewidth}
\begin{tcolorbox}[fmbase,colback=fmdefbg,colframe=fmdef,coltitle=white,title={Propiedad distributiva}]
{\centering$a\cdot(b\pm c)=a\cdot b\pm a\cdot c$\par}\smallskip
Para calcular de cabeza: \ $7\cdot99=7\cdot(100-1)=700-7=693$\\[3pt]
Para sacar \textbf{factor común}: \ $12\cdot15+12\cdot5=12\cdot20=240$
\end{tcolorbox}
\end{minipage}

\vfill
\begin{minipage}[c]{0.4\linewidth}\centering
\begin{tikzpicture}[font=\Large]
  \node (D) at (0,0) {$247$}; \node (d) at (1.6,0) {$12$};
  \draw[fmnavy,line width=1.2pt] (0.75,0.35) -- (0.75,-0.35) -- (2.4,-0.35);
  \node[text=fmA] at (1.6,-0.75) {$20$}; \node[text=fmR] at (0,-1.1) {$7$};
  \node[font=\scriptsize,text=fmmuted] at (-0.9,0) {dividendo}; \node[font=\scriptsize,text=fmmuted] at (2.7,0) {divisor};
  \node[font=\scriptsize,text=fmA] at (2.7,-0.75) {cociente}; \node[font=\scriptsize,text=fmR] at (-0.9,-1.1) {resto};
\end{tikzpicture}
\end{minipage}\hfill
\begin{minipage}[c]{0.57\linewidth}
\begin{idea}[La división entera]
Si no es exacta, da un \textbf{cociente} $q$ y un \textbf{resto} $r$, siempre \textbf{menor} que el divisor:\\[4pt]
{\centering\Large$D=d\cdot q+r$ \qquad \normalsize$247=12\cdot20+7$\par}\smallskip
Es la \textbf{prueba de la división}: si no se cumple, hay un error.
\end{idea}
\end{minipage}

\vfill
\begin{exemple}[Ejemplos]
\begin{minipage}[t]{0.48\linewidth}
250 caramelos en bolsas de 12: \ $250=12\cdot20+10$\\ \hspace*{1em}\boxed{20\text{ bolsas y sobran }10}\\[4pt]
$25\cdot101=25\cdot100+25=\boxed{2\,525}$
\end{minipage}\hfill
\begin{minipage}[t]{0.49\linewidth}
¿Qué dividendo tiene divisor 15, cociente 8 y resto 6?\\ \hspace*{1em}$D=15\cdot8+6=\boxed{126}$\\[4pt]
$36\cdot7+36\cdot3=36\cdot10=\boxed{360}$
\end{minipage}
\end{exemple}

\vfill
\begin{compte}[El resto siempre es menor que el divisor]
$100=7\cdot13+9$ \mal{} \ no es una división entera: $9>7$. \ La correcta: \ $100=7\cdot14+2$ \be{} \qquad Y \textbf{no se puede dividir entre 0}.
\end{compte}

\vfill
% ============================================================ 1.2
\clearpage
\seccio{sec2}{Potencias, raíces y jerarquía}{Multiplicaciones abreviadas y el orden en que operamos.}

\vfill
\begin{minipage}[c]{0.52\linewidth}
\begin{idea}[Potencia]
Una multiplicación de factores iguales: la \textbf{base} se repite tantas veces como dice el \textbf{exponente}.\\[4pt]
{\centering\Large$a^n=\underbrace{a\cdot a\cdots a}_{n\text{ veces}}$\par}\smallskip
$2^5=32$, \ $10^3=1000$. \ \ $a^2$: «al cuadrado»; $a^3$: «al cubo». \ $a^1=a$, \ $a^0=1$.
\end{idea}
\end{minipage}\hfill
\begin{minipage}[c]{0.45\linewidth}
\begin{compte}[El exponente no multiplica]
\be{2^5=2\cdot2\cdot2\cdot2\cdot2=32}\\[3pt]
\mal{2^5=2\cdot5=10}
\end{compte}
\end{minipage}

\vfill
\begin{idea}[Raíz cuadrada]
\begin{minipage}[c]{0.4\linewidth}
{\Large$\sqrt a=b\iff b^2=a$}\\[4pt]
$\sqrt{49}=7$ porque $7^2=49$.\\[3pt]
Si no es exacta, la \textbf{raíz entera}: \ $\sqrt{50}\approx7$, resto 1 \ ($50=7^2+1$).
\end{minipage}\hfill
\begin{minipage}[c]{0.57\linewidth}\centering\small
\taularod{>{\heavy}c*{6}{>{\hspace{3pt}}c<{\hspace{3pt}}}}{
\rowcolor{fmblue}\color{white}$n$ & \color{white}1 & \color{white}2 & \color{white}3 & \color{white}4 & \color{white}5 & \color{white}6\\
$n^2$ & 1 & 4 & 9 & 16 & 25 & 36\\
\rowcolor{fmblue}\color{white}$n$ & \color{white}7 & \color{white}8 & \color{white}9 & \color{white}10 & \color{white}11 & \color{white}12\\
$n^2$ & 49 & 64 & 81 & 100 & 121 & 144
}\\[3pt]
{\comic\color{fmmuted}los cuadrados perfectos: su raíz es exacta}
\end{minipage}
\end{idea}

\vfill
\begin{minipage}[c]{0.42\linewidth}\centering
\begin{tikzpicture}[st/.style={draw=#1,fill=#1!12,line width=1.2pt,rounded corners=2mm,minimum height=0.8cm,anchor=south west,font=\large,inner xsep=6pt}]
  \node[st=fmV] at (0,2.7) {\pas{1}\ \ Paréntesis $(\ )\ [\ ]$};
  \node[st=fmA] at (0.5,1.8) {\pas{2}\ \ Potencias y raíces};
  \node[st=fmK] at (1.0,0.9) {\pas{3}\ \ $\cdot$ \ y \ $:$};
  \node[st=fmB] at (1.5,0) {\pas{4}\ \ $+$ \ y \ $-$};
\end{tikzpicture}\\[3pt]
{\comic\small\color{fmmuted}\pas{3} y \pas{4}: de izquierda a derecha}
\end{minipage}\hfill
\begin{minipage}[c]{0.55\linewidth}
\begin{exemple}[Ejemplo: \ $20-3^2\cdot2+\sqrt{16}$]
$=20-9\cdot2+4$ \ \nota{\pas{2}}\\[3pt]
$=20-18+4$ \ \nota{\pas{3}}\\[3pt]
$=\boxed{6}$ \ \nota{\pas{4}}
\end{exemple}
\end{minipage}

\vfill
\begin{compte}[Errores frecuentes]
\begin{tabular}{@{}p{0.31\linewidth}p{0.31\linewidth}p{0.3\linewidth}@{}}
\be{8+2\cdot5=18}\newline{\small primero el producto, no $10\cdot5$} &
\be{12:2\cdot3=18}\newline{\small de izquierda a derecha, no $12:6$} &
\be{\sqrt{9+16}=\sqrt{25}=5}\newline{\small no $\sqrt9+\sqrt{16}=7$}
\end{tabular}
\end{compte}

\vfill
\begin{exemple}[Ejemplos]
\begin{minipage}[t]{0.48\linewidth}
$(3+4)\cdot5-2^3=35-8=\boxed{27}$\\[4pt]
$\sqrt{49}+2\cdot(8-3)^2=7+2\cdot25=\boxed{57}$
\end{minipage}\hfill
\begin{minipage}[t]{0.49\linewidth}
$50-[18:(3^2-6)]\cdot4=50-[18:3]\cdot4=50-24=\boxed{26}$\\[4pt]
$5\cdot\sqrt{36}-2^4:8=30-16:8=30-2=\boxed{28}$
\end{minipage}
\end{exemple}

\vfill
% ============================================================ 1.3
\clearpage
\seccio{sec3}{Valor absoluto, opuesto y orden}{En la recta, cuanto más a la derecha, mayor.}

\vfill
\begin{center}
\begin{tikzpicture}[x=0.95cm]
  \recta{-7}{7}
  \fill[fmR] (-5,0) circle (4pt); \fill[fmR] (5,0) circle (4pt);
  \draw[decorate,decoration={brace,amplitude=6pt},fmA,line width=1.2pt] (-5,0.35) -- (0,0.35) node[midway,above=7pt,font=\large,text=fmA]{$|-5|=5$};
  \draw[decorate,decoration={brace,amplitude=6pt},fmA,line width=1.2pt] (0,0.35) -- (5,0.35) node[midway,above=7pt,font=\large,text=fmA]{$|5|=5$};
  \draw[{Latex[length=2.5mm]}-{Latex[length=2.5mm]},fmB,line width=1.3pt] (-5,-0.75) to[bend right=18] node[below,font=\large,text=fmB]{opuestos: \ $\text{op}(5)=-5$} (5,-0.75);
  \node[font=\comic,text=fmmuted,anchor=east] at (-6.6,0.8) {menores};
  \node[font=\comic,text=fmmuted,anchor=west] at (6.6,0.8) {mayores};
\end{tikzpicture}
\end{center}

\vfill
\begin{tcbraster}[raster columns=2,raster equal height,raster column skip=10pt]
\begin{tcolorbox}[fmbase,colback=fmdefbg,colframe=fmdef,coltitle=white,title={Valor absoluto \ $|a|$}]
La \textbf{distancia} del número al 0. Nunca es negativo.\par\smallskip
{\centering$|-7|=7\qquad |7|=7\qquad |0|=0$\par}
\end{tcolorbox}
\begin{tcolorbox}[fmbase,colback=fmdefbg,colframe=fmdef,coltitle=white,title={Opuesto \ $\text{op}(a)=-a$}]
Mismo valor absoluto, \textbf{signo contrario}.\par\smallskip
{\centering$\text{op}(5)=-5\qquad \text{op}(-3)=3\qquad a+\text{op}(a)=0$\par}
\end{tcolorbox}
\end{tcbraster}

\vfill
\subsection{Ordenar enteros}
\begin{minipage}[t]{0.52\linewidth}
\begin{idea}
\begin{itemize}[leftmargin=1em]
\item Cualquier \textbf{positivo} es mayor que 0, y 0 es mayor que cualquier \textbf{negativo}.
\item Entre dos \textbf{negativos}, es mayor el que tiene \textbf{menor} valor absoluto: \ $-3>-8$.
\end{itemize}
\end{idea}
\end{minipage}\hfill
\begin{minipage}[t]{0.45\linewidth}
\begin{truc}[Piensa en el termómetro]
Hace menos frío a $-3$\,°C que a $-8$\,°C:\\[3pt]
\hspace*{1em}{\Large$-3>-8$} \ \ {\color{fmA}\Large\faThermometerQuarter}
\end{truc}
\end{minipage}

\vfill
\begin{exemple}[Ejemplos]
\begin{minipage}[t]{0.5\linewidth}
Ordena de menor a mayor: \ $-5,\ 3,\ -1,\ 0,\ -12,\ 7$\\[4pt]
{\large$-12<-5<-1<0<3<7$}
\end{minipage}\hfill
\begin{minipage}[t]{0.47\linewidth}
$|-4|+|3|=4+3=7$\\[3pt]
$\text{op}\big(|-6|\big)=\text{op}(6)=-6$\\[3pt]
$-|-2|=-2$
\end{minipage}
\end{exemple}

\vfill
\begin{exemple}[Ejemplo: \ ¿qué enteros cumplen \ $|x|=5$? \ ¿Y \ $|x|<3$?]
\begin{minipage}[c]{0.38\linewidth}
$|x|=5$: \ $\boxed{x=5}$ \ o \ $\boxed{x=-5}$\\[6pt]
$|x|<3$: \ $\boxed{-2,\ -1,\ 0,\ 1,\ 2}$\\[3pt]
\nota{a menos de 3 pasos del 0}
\end{minipage}\hfill
\begin{minipage}[c]{0.6\linewidth}\centering
\begin{tikzpicture}[x=0.7cm]
  \recta{-5}{5}
  \foreach \k in {-2,...,2} \fill[fmK] (\k,0) circle (4.5pt);
  \fill[fmB] (-5,0) circle (4.5pt); \fill[fmB] (5,0) circle (4.5pt);
  \draw[decorate,decoration={brace,amplitude=5pt},fmK,line width=1.1pt] (-2,0.3) -- (2,0.3) node[midway,above=6pt,font=\small,text=fmK!80!black]{$|x|<3$};
\end{tikzpicture}
\end{minipage}
\end{exemple}

\vfill
% ============================================================ 1.4
\clearpage
\seccio{sec4}{Operaciones combinadas con enteros}{Sumas, productos y la jerarquía de siempre.}

\vfill
\begin{tcbraster}[raster columns=2,raster equal height,raster column skip=10pt]
\begin{tcolorbox}[fmbase,colback=fmK!8,colframe=fmK,coltitle=white,title={Mismo signo: \textbf{sumamos}}]
Sumamos los valores absolutos y ponemos el \textbf{signo común}.\par\smallskip
{\centering$(+4)+(+3)=+7\qquad (-4)+(-3)=-7$\par}
\end{tcolorbox}
\begin{tcolorbox}[fmbase,colback=fmB!8,colframe=fmB,coltitle=white,title={Signos distintos: \textbf{restamos}}]
El mayor menos el menor, con el signo del que tiene \textbf{mayor} valor absoluto.\par\smallskip
{\centering$(+4)+(-9)=-5\qquad (-4)+(+9)=+5$\par}
\end{tcolorbox}
\end{tcbraster}

\vfill
\begin{minipage}[c]{0.46\linewidth}
\begin{idea}[Quitar paréntesis]
{\centering\renewcommand{\arraystretch}{1.4}\Large
\begin{tabular}{cccc}
$+(+a)$ & $+(-a)$ & $-(+a)$ & $-(-a)$\\
$\cK{+a}$ & $\cR{-a}$ & $\cR{-a}$ & $\cK{+a}$
\end{tabular}\par}\smallskip
Un $-$ delante del paréntesis cambia \textbf{todos} los signos: \ $-(3-5+2)=-3+5-2$
\end{idea}
\end{minipage}\hfill
\begin{minipage}[c]{0.51\linewidth}
\begin{exemple}[Ejemplo: \ $7-(-3)+(-5)-(+2)$]
$=7+3-5-2$ \ \nota{quitamos paréntesis}\\[3pt]
$=10-7=\boxed{3}$ \ \nota{positivos y negativos}
\end{exemple}
\begin{truc}[Agrupa]
Suma aparte los positivos y los negativos:\\[2pt]
$-4+9-12+3-7=(9+3)-(4+12+7)=12-23=-11$
\end{truc}
\end{minipage}

\vfill
\begin{minipage}[t]{0.48\linewidth}
\begin{exemple}[Ejemplo: \ $-8+(4-9)-(-6+2)$]
$=-8+(-5)-(-4)$ \ \nota{dentro de los paréntesis}\\[3pt]
$=-8-5+4$ \ \nota{quitamos paréntesis}\\[3pt]
$=-13+4=\boxed{-9}$
\end{exemple}
\end{minipage}\hfill
\begin{minipage}[t]{0.49\linewidth}
\begin{exemple}[{Ejemplo: \ $15-\left[6-(2-9)\right]$}]
$=15-\left[6-(-7)\right]$ \ \nota{el de dentro primero}\\[3pt]
$=15-\left[6+7\right]$\\[3pt]
$=15-13=\boxed{2}$
\end{exemple}
\end{minipage}

\vfill
\begin{compte}[El signo menos delante del paréntesis]
\be{5-(-3)=5+3=8}\qquad\quad \be{-2-6=-2+(-6)=-8}\qquad\quad \mal{-(4-9)=-4-9}
\end{compte}

\vfill
\clearpage
\subsection{Multiplicación y división de enteros}

\vfill
\begin{minipage}[c]{0.36\linewidth}\centering
\renewcommand{\arraystretch}{1.6}\taulaclip{fmblue,line width=1.2pt}{\fontsize{20}{24}\selectfont%
\begin{tabular}{>{\columncolor{fmblue}\color{white}}w{c}{1.6cm}w{c}{1.4cm}w{c}{1.4cm}}
\rowcolor{fmblue}{\large$\cdot$ \ $:$} & \color{white}$+$ & \color{white}$-$\\
$+$ & \cellcolor{fmK!18}$\cK{+}$ & \cellcolor{fmR!14}$\cR{-}$\\
$-$ & \cellcolor{fmR!14}$\cR{-}$ & \cellcolor{fmK!18}$\cK{+}$
\end{tabular}}
\end{minipage}\hfill
\begin{minipage}[c]{0.6\linewidth}
\begin{idea}[La regla de los signos]
{\large Multiplicamos (o dividimos) los valores absolutos y ponemos el signo de la tabla.}\\[8pt]
$(-6)\cdot(+4)=\cR{-}24\qquad (-3)\cdot(-5)=\cK{+}15$\\[4pt]
$(-36):(-4)=\cK{+}9\qquad 42:(-7)=\cR{-}6$
\end{idea}
\end{minipage}

\vfill
\begin{minipage}[t]{0.5\linewidth}
\begin{idea}[Cuenta los negativos]
\textbf{Par} de negativos $\to$ \cK{positivo}. \ \textbf{Impar} $\to$ \cR{negativo}.\\[4pt]
$(-2)\cdot3\cdot(-1)\cdot(-5)=-30$ \ \nota{3 negativos}\\[3pt]
$(-1)^{15}=-1\qquad (-1)^{20}=+1$
\end{idea}
\end{minipage}\hfill
\begin{minipage}[t]{0.47\linewidth}
\begin{truc}[¿Y si no es exacta?]
La división solo da un entero si es \textbf{exacta}. Si no, da una \textbf{fracción}:\\[3pt]
$(-7):2=-\dfrac72$ \ \ \nota{¡por eso hacen falta los racionales!}
\end{truc}
\end{minipage}

\vfill
\begin{exemple}[Ejemplos]
\begin{minipage}[t]{0.48\linewidth}
$(-1)\cdot(-2)\cdot(-3)=\cR{-6}$ \ \nota{3 negativos}\\[4pt]
$(-1)\cdot(-2)\cdot(-3)\cdot(-4)=\cK{+24}$ \ \nota{4 negativos}
\end{minipage}\hfill
\begin{minipage}[t]{0.49\linewidth}
$(-60):(+5):(-3)=$ \ \nota{de izquierda a derecha}\\[4pt]
$=(-12):(-3)=\boxed{4}$
\end{minipage}
\end{exemple}

\vfill
\begin{compte}[Divisiones seguidas: de izquierda a derecha]
\be{(-100):(-4):(-5)=25:(-5)=-5}\qquad\quad \mal{(-100):\big((-4):(-5)\big)}
\end{compte}

\vfill
\begin{tcbraster}[raster columns=2,raster equal height,raster column skip=10pt]
\begin{tcolorbox}[fmbase,colback=fmdefbg,colframe=fmdef,coltitle=white,title={\faSuperscript\ \ Potencias de base negativa}]
Exponente \textbf{par} $\to$ \cK{positivo}. \ \textbf{Impar} $\to$ \cR{negativo}.\par\smallskip
{\centering$(-2)^4=\cK{+16}\qquad (-2)^3=\cR{-8}$\par}
\end{tcolorbox}
\begin{compte}[¿Con paréntesis o sin?]
\be{(-2)^4=(-2)\cdot(-2)\cdot(-2)\cdot(-2)=16}\\[3pt]
\be{-2^4=-(2\cdot2\cdot2\cdot2)=-16}\\[2pt]
{\small Sin paréntesis, el $-$ no se eleva.}
\end{compte}
\end{tcbraster}

% ------------------------------------------------------------
\clearpage

\vfill
\clearpage
\subsection{Operaciones combinadas con enteros}

\vfill
\begin{minipage}[c]{0.42\linewidth}\centering
\begin{tikzpicture}[st/.style={draw=#1,fill=#1!12,line width=1.2pt,rounded corners=2mm,minimum height=0.8cm,anchor=south west,font=\large,inner xsep=6pt}]
  \node[st=fmV] at (0,2.7) {\pas{1}\ \ Paréntesis $(\ )\ [\ ]$};
  \node[st=fmA] at (0.5,1.8) {\pas{2}\ \ Potencias};
  \node[st=fmK] at (1.0,0.9) {\pas{3}\ \ $\cdot$ \ y \ $:$};
  \node[st=fmB] at (1.5,0) {\pas{4}\ \ $+$ \ y \ $-$};
\end{tikzpicture}\\[3pt]
{\comic\small\color{fmmuted}\pas{3} y \pas{4}: de izquierda a derecha}
\end{minipage}\hfill
\begin{minipage}[c]{0.55\linewidth}
\begin{exemple}[Ejemplo: \ $12-3\cdot(5-8)+20:(-4)$]
\vspace{-6pt}
\begin{alignat*}{2}
&=12-3\cdot(-3)+20:(-4) &\quad& \text{\nota{\pas{1}}}\\
&=12-(-9)+(-5) && \text{\nota{\pas{3}}}\\
&=12+9-5=\boxed{16} && \text{\nota{\pas{4}}}
\end{alignat*}
\end{exemple}
\end{minipage}

\vfill
\begin{compte}[Errores típicos]
\begin{tabular}{@{}p{0.31\linewidth}p{0.31\linewidth}p{0.3\linewidth}@{}}
\be{12:2\cdot3=6\cdot3=18}\newline{\small de izquierda a derecha, no $12:6$} &
\be{5-2\cdot(-3)=5+6=11}\newline{\small primero el producto, no $3\cdot(-3)$} &
\be{-3^2=-9}\newline\be{(-3)^2=9}
\end{tabular}
\end{compte}

\vfill
\begin{minipage}[t]{0.48\linewidth}
\begin{exemple}[{Ejemplo: \ $(-2)^3-4\cdot\left[5-(3-6)^2\right]$}]
$=(-2)^3-4\cdot\left[5-(-3)^2\right]$ \ \nota{\pas{1}}\\[3pt]
$=(-2)^3-4\cdot[5-9]$ \ \nota{\pas{2}}\\[3pt]
$=(-2)^3-4\cdot(-4)$\\[3pt]
$=-8+16=\boxed{8}$ \ \nota{\pas{2} \pas{3} \pas{4}}
\end{exemple}
\end{minipage}\hfill
\begin{minipage}[t]{0.49\linewidth}
\begin{exemple}[Ejemplo: \ $18:(-3)\cdot2-(-5)^2+3\cdot(-2)^2$]
$=18:(-3)\cdot2-25+3\cdot4$ \ \nota{\pas{2}}\\[3pt]
$=(-6)\cdot2-25+12$ \ \nota{\pas{3}}\\[3pt]
$=-12-25+12$\\[3pt]
$=\boxed{-25}$ \ \nota{\pas{4}}
\end{exemple}
\end{minipage}

\vfill
\begin{minipage}[t]{0.48\linewidth}
\begin{exemple}[Ejemplo: \ $3-2\cdot(4-7)-(-6):2$]
$=3-2\cdot(-3)-(-6):2$ \ \nota{\pas{1}}\\[3pt]
$=3-(-6)-(-3)$ \ \nota{\pas{3}}\\[3pt]
$=3+6+3=\boxed{12}$ \ \nota{\pas{4}}
\end{exemple}
\end{minipage}\hfill
\begin{minipage}[t]{0.49\linewidth}
\begin{exemple}[{Ejemplo: \ $2\cdot\left\{3-\left[4-(1-6)\right]\right\}$}]
$=2\cdot\left\{3-\left[4-(-5)\right]\right\}$ \ \nota{de dentro hacia fuera}\\[3pt]
$=2\cdot\left\{3-9\right\}$\\[3pt]
$=2\cdot(-6)=\boxed{-12}$
\end{exemple}
\end{minipage}

\vfill
% ============================================================ 1.5
\clearpage
\seccio{sec5}{Fracciones y números racionales}{Partimos la unidad en partes iguales y tomamos unas cuantas.}

\vfill
\begin{minipage}[c]{0.45\linewidth}\centering
\begin{tikzpicture}
  \node[font=\fontsize{40}{40}\selectfont] (f) at (0,0) {$\dfrac{\cB{3}}{\cA{4}}$};
  \node[anchor=west,font=\large,text=fmB,align=left] (n) at (1.4,0.75) {\textbf{numerador}: tomamos 3};
  \node[anchor=west,font=\large,text=fmA,align=left] (d) at (1.4,-0.75) {\textbf{denominador}: 4 partes};
  \draw[-{Latex[length=2.5mm]},fmB,line width=1.2pt] (n.west) -- ([xshift=5pt,yshift=13pt]f.east);
  \draw[-{Latex[length=2.5mm]},fmA,line width=1.2pt] (d.west) -- ([xshift=5pt,yshift=-13pt]f.east);
\end{tikzpicture}\\[6pt]
\barra{4}{3}{fmB}{5}\\[2pt]
{\comic\color{fmmuted}también es la división \ $3:4=0{,}75$}
\end{minipage}\hfill
\begin{minipage}[c]{0.52\linewidth}
\begin{idea}[Fracciones equivalentes]
Representan el \textbf{mismo} número. Lo comprobamos \textbf{en cruz}:\\[4pt]
{\centering
\begin{tikzpicture}[baseline=(eq.base),n/.style={font=\Large,inner sep=2pt}]
  \node[n,text=fmB] (a) at (0,0.42) {$6$}; \node[n,text=fmA] (b) at (0,-0.42) {$8$};
  \draw[line width=0.8pt] (-0.3,0) -- (0.3,0);
  \node[n] (eq) at (0.85,0) {$=$};
  \node[n,text=fmA] (c) at (1.7,0.42) {$15$}; \node[n,text=fmB] (d) at (1.7,-0.42) {$20$};
  \draw[line width=0.8pt] (1.35,0) -- (2.05,0);
  \draw[fmB,line width=1.2pt,-{Latex[length=2mm]},shorten <=1pt,shorten >=1pt] (a.east) -- (d.west);
  \draw[fmA,line width=1.2pt,-{Latex[length=2mm]},shorten <=1pt,shorten >=1pt] (b.east) -- (c.west);
\end{tikzpicture}\qquad $\cB{6\cdot20}=\cA{8\cdot15}=120$\par}
\smallskip Un \textbf{número racional} son todas las fracciones equivalentes; lo escribimos con la \textbf{irreducible}.
\end{idea}
\end{minipage}

\vfill
\begin{tcbraster}[raster columns=2,raster equal height,raster column skip=10pt]
\begin{tcolorbox}[fmbase,colback=fmdefbg,colframe=fmdef,coltitle=white,title={\faSearchPlus\ \ Amplificar}]
Multiplicamos arriba y abajo por el mismo número:\par\smallskip
{\centering$\dfrac23=\dfrac{2\cdot\cK{4}}{3\cdot\cK{4}}=\dfrac{8}{12}$\par}
\end{tcolorbox}
\begin{tcolorbox}[fmbase,colback=fmdefbg,colframe=fmdef,coltitle=white,title={\faSearchMinus\ \ Simplificar}]
Dividimos por el mismo número (por el \textbf{m.c.d.}, de una vez):\par\smallskip
{\centering$\dfrac{18}{24}=\dfrac{18:\cK{6}}{24:\cK{6}}=\dfrac34$\par}
\end{tcolorbox}
\end{tcbraster}

\vfill
\begin{minipage}[t]{0.5\linewidth}
\begin{idea}[Comparar fracciones]
Las ponemos con el \textbf{mismo denominador} (el m.c.m.) y miramos los numeradores:\\[4pt]
$\dfrac56$ y $\dfrac79$: \ m.c.m.$(6,9)=18$\\[3pt]
$\dfrac{15}{18}>\dfrac{14}{18}\ \to\ \dfrac56>\dfrac79$
\end{idea}
\end{minipage}\hfill
\begin{minipage}[t]{0.47\linewidth}
\begin{truc}[El signo]
Lo pongas donde lo pongas, vale igual:\\[3pt]
$-\dfrac34=\dfrac{-3}{4}=\dfrac{3}{-4}$\\[3pt]
Y menos entre menos: \ $\dfrac{-3}{-4}=\dfrac34$
\end{truc}
\end{minipage}

\vfill
\begin{exemple}[Ejemplo: \ simplifica \ $\dfrac{84}{126}$]
$84=2^2\cdot3\cdot7$, \ $126=2\cdot3^2\cdot7$ \ $\to$ \ m.c.d.$=2\cdot3\cdot7=42$ \ $\to$ \ $\dfrac{84:42}{126:42}=\boxed{\dfrac23}$
\end{exemple}

\vfill
\begin{exemple}[Ejemplo: \ ordena de menor a mayor \ $-\dfrac12,\ \ \dfrac34,\ \ -\dfrac23,\ \ \dfrac56$]
\begin{minipage}[c]{0.5\linewidth}
m.c.m.$(2,4,3,6)=12$:\\[4pt]
$-\dfrac{6}{12},\ \ \dfrac{9}{12},\ \ -\dfrac{8}{12},\ \ \dfrac{10}{12}$
\end{minipage}\hfill
\begin{minipage}[c]{0.47\linewidth}
Comparamos numeradores: \ $-8<-6<9<10$\\[4pt]
{\large$-\dfrac23<-\dfrac12<\dfrac34<\dfrac56$}
\end{minipage}
\end{exemple}

\vfill
% ============================================================ 1.6
\clearpage
\seccio{sec6}{Operaciones con fracciones}{Sumar pide el mismo denominador; multiplicar, no.}

\vfill
\begin{tcbraster}[raster columns=2,raster equal height=rows,raster column skip=10pt]
\begin{tcolorbox}[fmbase,colback=fmdefbg,colframe=fmdef,coltitle=white,title={\faPlus\,\faMinus\ \ Suma y resta}]
Con el \textbf{mismo denominador}, sumamos los numeradores. Si no, primero el \fmkey{m.c.m.}\par\smallskip
{\centering$\dfrac56-\dfrac34=\dfrac{10}{12}-\dfrac{9}{12}=\dfrac{1}{12}$\par}
\end{tcolorbox}
\begin{tcolorbox}[fmbase,colback=fmdefbg,colframe=fmdef,coltitle=white,title={\faTimes\ \ Multiplicación}]
\textbf{Arriba por arriba} y \textbf{abajo por abajo}.\par\smallskip
{\centering$\dfrac{a}{b}\cdot\dfrac{c}{d}=\dfrac{a\cdot c}{b\cdot d}\qquad \dfrac23\cdot\dfrac45=\dfrac{8}{15}$\par}
\end{tcolorbox}
\begin{tcolorbox}[fmbase,colback=fmdefbg,colframe=fmdef,coltitle=white,title={\faDivide\ \ División}]
Multiplicamos \textbf{en cruz} (o por la \textbf{inversa}).\par\smallskip
{\centering$\dfrac{a}{b}:\dfrac{c}{d}=\dfrac{a\cdot d}{b\cdot c}\qquad \dfrac23:\dfrac45=\dfrac{10}{12}=\dfrac56$\par}
\end{tcolorbox}
\begin{tcolorbox}[fmbase,colback=fmdefbg,colframe=fmdef,coltitle=white,title={\faUndo\ \ Inversa}]
La inversa de $\dfrac{a}{b}$ es $\dfrac{b}{a}$: su producto es 1.\par\smallskip
{\centering La de $-2$ es $-\dfrac12$. \ El 0 no tiene.\par}
\end{tcolorbox}
\end{tcbraster}

\vfill
\begin{minipage}[t]{0.47\linewidth}
\begin{truc}[Un entero es una fracción]
Tiene denominador 1:\\[3pt]
$3-\dfrac75=\dfrac{15}{5}-\dfrac75=\dfrac85$
\end{truc}
\end{minipage}\hfill
\begin{minipage}[t]{0.5\linewidth}
\begin{truc}[Simplifica antes de multiplicar]
Descompón y tacha los factores comunes:\\[3pt]
$-\dfrac{4\cdot15}{9\cdot8}=-\dfrac{\cR{2^2}\cdot\cR{3}\cdot5}{\cR{3}\cdot3\cdot\cR{2^2}\cdot2}=-\dfrac{5}{6}$
\end{truc}
\end{minipage}

\vfill
\begin{exemple}[Ejemplos]
\begin{minipage}[t]{0.32\linewidth}
$\dfrac56-\dfrac34+\dfrac12=$\\[3pt]
\nota{m.c.m.$(6,4,2)=12$}\\[2pt]
$=\dfrac{10-9+6}{12}=\boxed{\dfrac{7}{12}}$
\end{minipage}\hfill
\begin{minipage}[t]{0.32\linewidth}
$\left(-\dfrac49\right)\cdot\dfrac{15}{8}=$\\[3pt]
\nota{signo: $-\cdot+=-$}\\[2pt]
$=-\dfrac{60}{72}=\boxed{-\dfrac56}$
\end{minipage}\hfill
\begin{minipage}[t]{0.32\linewidth}
$\dfrac{6}{25}:\left(-\dfrac{9}{10}\right)=$\\[3pt]
\nota{en cruz}\\[2pt]
$=-\dfrac{60}{225}=\boxed{-\dfrac{4}{15}}$
\end{minipage}
\end{exemple}

\vfill
\begin{compte}[Errores típicos]
\begin{tabular}{@{}p{0.47\linewidth}p{0.47\linewidth}@{}}
\mal{\dfrac12+\dfrac13=\dfrac25}\ \ {\small no sumes los denominadores}\par\vspace{9pt}\be{\dfrac12+\dfrac13=\dfrac36+\dfrac26=\dfrac56} &
\mal{\dfrac23\cdot\dfrac45=\dfrac{10}{12}}\ \ {\small eso es dividir}\par\vspace{9pt}\be{\dfrac23\cdot\dfrac45=\dfrac{8}{15}}
\end{tabular}
\end{compte}

\vfill
% ============================================================ 1.7
\clearpage
\seccio{sec7}{Operaciones combinadas con fracciones}{La misma jerarquía que con los enteros.}

\vfill
\begin{minipage}[c]{0.4\linewidth}\raggedright
\pas{1} Paréntesis $\to$ \pas{2} potencias $\to$ \pas{3} $\cdot$ y $:$ $\to$ \pas{4} $+$ y $-$.\\[8pt]
\faCheck\ \textbf{Simplifica} cada resultado parcial.\\[6pt]
\faCheck\ Común denominador \textbf{solo} para los términos que sumas en ese paso.\\[6pt]
\faCheck\ Un $-$ delante de un paréntesis afecta a \textbf{todo} el resultado.
\end{minipage}\hfill
\begin{minipage}[c]{0.57\linewidth}
\begin{exemple}[Ejemplo: \ $\dfrac34-\dfrac12\cdot\left(\dfrac56-\dfrac13\right)$]
\vspace{-6pt}
\begin{alignat*}{2}
&=\dfrac34-\dfrac12\cdot\left(\dfrac56-\dfrac26\right)=\dfrac34-\dfrac12\cdot\dfrac12 &\quad& \text{\nota{\pas{1}}}\\
&=\dfrac34-\dfrac14 && \text{\nota{\pas{3}}}\\
&=\dfrac24=\boxed{\dfrac12} && \text{\nota{\pas{4}}}
\end{alignat*}
\end{exemple}
\end{minipage}

\vfill
\begin{exemple}[Ejemplo: \ $\left(1-\dfrac25\right)^{2}:\dfrac{9}{10}+\dfrac13$]
$=\left(\dfrac35\right)^{2}:\dfrac{9}{10}+\dfrac13=\dfrac{9}{25}:\dfrac{9}{10}+\dfrac13=\dfrac{90}{225}+\dfrac13=\dfrac25+\dfrac13=\dfrac{6}{15}+\dfrac{5}{15}=\boxed{\dfrac{11}{15}}$
\end{exemple}

\vfill
\begin{minipage}[t]{0.48\linewidth}
\begin{exemple}[{Ejemplo: \ $2-\left[\dfrac13-\left(\dfrac12-\dfrac34\right)\right]$}]
$=2-\left[\dfrac13-\left(-\dfrac14\right)\right]$ \ \nota{de dentro hacia fuera}\\[3pt]
$=2-\left[\dfrac4{12}+\dfrac3{12}\right]=2-\dfrac7{12}=\boxed{\dfrac{17}{12}}$
\end{exemple}
\end{minipage}\hfill
\begin{minipage}[t]{0.49\linewidth}
\begin{exemple}[Ejemplo: \ $\dfrac23:\dfrac49-\dfrac56\cdot\dfrac35$]
$=\dfrac{18}{12}-\dfrac{15}{30}$ \ \nota{\pas{3} primero}\\[3pt]
$=\dfrac32-\dfrac12=\boxed{1}$
\end{exemple}
\end{minipage}

\vfill
\begin{exemple}[Ejemplo: \ $\left(-\dfrac12\right)^2\cdot\left(3-\dfrac13\right)+\dfrac56:\left(-\dfrac53\right)$]
$=\dfrac14\cdot\dfrac83+\dfrac56\cdot\left(-\dfrac35\right)=\dfrac{8}{12}-\dfrac{15}{30}=\dfrac23-\dfrac12=\boxed{\dfrac16}$ \ \nota{dividir es multiplicar por la inversa}
\end{exemple}

\vfill
\begin{compte}[El signo menos delante de una fracción]
$\dfrac12-\left(\dfrac13-\dfrac14\right)=\dfrac12-\dfrac1{12}=\dfrac5{12}$ \qquad \mal{\dfrac12-\dfrac13-\dfrac14}: el $-$ cambia \textbf{todos} los signos del paréntesis.
\end{compte}

\vfill
% ============================================================ 1.8
\clearpage
\seccio{sec8}{Problemas con fracciones}{Un dibujo con una barra lo aclara todo.}

\vfill
\begin{tcbraster}[raster columns=2,raster equal height=rows,raster column skip=10pt]
\begin{tcolorbox}[fmbase,colback=fmdefbg,colframe=fmdef,coltitle=white,title={Fracción de una cantidad}]
Dividimos entre el denominador y multiplicamos por el numerador.\par\smallskip
{\centering$\dfrac34$ de $200=\dfrac{3\cdot200}{4}=150$\par}
\end{tcolorbox}
\begin{tcolorbox}[fmbase,colback=fmdefbg,colframe=fmdef,coltitle=white,title={Fracción de una fracción}]
Las \textbf{multiplicamos}.\par\smallskip
{\centering$\dfrac23$ de $\dfrac35=\dfrac23\cdot\dfrac35=\dfrac25$\par}
\end{tcolorbox}
\begin{tcolorbox}[fmbase,colback=fmdefbg,colframe=fmdef,coltitle=white,title={La parte que queda}]
El total es $1=\dfrac{b}{b}$. Si gastamos $\dfrac{a}{b}$, queda $1-\dfrac{a}{b}$.\par\smallskip
{\centering Gastamos $\dfrac25$: \ queda $\dfrac55-\dfrac25=\dfrac35$\par}
\end{tcolorbox}
\begin{tcolorbox}[fmbase,colback=fmdefbg,colframe=fmdef,coltitle=white,title={Problema inverso}]
Si $\dfrac{a}{b}$ del total es $P$, el total es $P:\dfrac{a}{b}$.\par\smallskip
{\centering$\dfrac23$ del total son 40: \ $40:\dfrac23=60$\par}
\end{tcolorbox}
\end{tcbraster}

\vfill
\begin{exemple}[Ejemplo: \ Ana gasta $\frac25$ de sus 60 € en ropa y $\frac13$ de lo que le queda en un libro. ¿Cuánto le queda?]
\begin{minipage}[c]{0.45\linewidth}\centering
\begin{tikzpicture}[x=1.3cm]
  \foreach \i in {0,...,4} {\filldraw[fill=\ifnum\i<2 fmB!35\else white\fi,draw=fmB,line width=1pt] (\i,0) rectangle (\i+1,0.6);}
  \node[font=\small] at (1,0.3) {ropa};
  \filldraw[fill=fmA!35,draw=fmA,line width=1pt] (2,0) rectangle (3,0.6); \node[font=\small] at (2.5,0.3) {libro};
  \node[font=\small] at (4,0.3) {queda};
  \draw[decorate,decoration={brace,amplitude=5pt},line width=1pt] (0,0.7) -- (5,0.7) node[midway,above=6pt]{60 €};
  \foreach \i in {0,...,4} \node[font=\small,text=fmmuted] at (\i+0.5,-0.3) {12 €};
\end{tikzpicture}
\end{minipage}\hfill
\begin{minipage}[c]{0.52\linewidth}
Ropa: \ $\dfrac25$ de $60=24$ € \ $\to$ \ quedan 36 €\\[4pt]
Libro: \ $\dfrac13$ de $36=12$ €\\[4pt]
Le quedan \ $36-12=\boxed{24\text{ €}}$
\end{minipage}
\end{exemple}

\vfill
\begin{exemple}[Ejemplo: \ un depósito está lleno a $\frac23$ y tiene 40 litros. ¿Qué capacidad tiene?]
\begin{minipage}[c]{0.4\linewidth}\centering
\begin{tikzpicture}[x=1.4cm]
  \foreach \i in {0,1,2} {\filldraw[fill=\ifnum\i<2 fmA!35\else white\fi,draw=fmA,line width=1pt] (\i,0) rectangle (\i+1,0.6);}
  \draw[decorate,decoration={brace,amplitude=5pt},line width=1pt] (0,0.7) -- (2,0.7) node[midway,above=6pt]{40 L};
  \foreach \i in {0,1,2} \node[font=\small,text=fmmuted] at (\i+0.5,-0.3) {20 L};
\end{tikzpicture}
\end{minipage}\hfill
\begin{minipage}[c]{0.57\linewidth}
Problema \textbf{inverso}: \ $40:\dfrac23=\dfrac{40\cdot3}{2}=\boxed{60\text{ L}}$\\[4pt]
\nota{Cada tercio son 20 L; los tres tercios, 60 L.}
\end{minipage}
\end{exemple}

\vfill
\begin{minipage}[t]{0.5\linewidth}
\begin{compte}[¿De qué cantidad?]
El $\frac13$ es \textbf{de lo que queda} (36 €), no del total (60 €). Lee bien de qué es cada fracción.
\end{compte}
\end{minipage}\hfill
\begin{minipage}[t]{0.47\linewidth}
\begin{truc}[Comprueba que tiene sentido]
Responde con una \textbf{frase} y con las \textbf{unidades}: «Le quedan 24 €.»
\end{truc}
\end{minipage}

\vfill
% ============================================================ 1.9
\clearpage
\seccio{sec9}{Potencias de base entera y racional}{Una potencia es una multiplicación de factores iguales.}

\vfill
\begin{center}
\begin{tikzpicture}
  \node[font=\fontsize{44}{44}\selectfont,inner sep=2pt] (p) {$\cA{3}^{\cB{4}}$};
  \node[font=\fontsize{26}{26}\selectfont,anchor=west] (e) at ([xshift=4pt]p.east) {$=\cA{3}\cdot\cA{3}\cdot\cA{3}\cdot\cA{3}=81$};
  \draw[decorate,decoration={brace,amplitude=6pt,mirror},fmB,line width=1.2pt] ([xshift=22pt,yshift=-6pt]e.south west) -- ([xshift=-58pt,yshift=-6pt]e.south east) node[midway,below=8pt,font=\large,text=fmB]{4 factores};
  \coordinate (pb) at ([xshift=-12pt,yshift=-16pt]p.center);
  \coordinate (pe) at ([xshift=10pt,yshift=31pt]p.center);
  \node[fill=fmA!14,draw=fmA,rounded corners=2mm,font=\large,anchor=east] (lb) at ([xshift=-1.6cm,yshift=-0.7cm]p.west) {{\heavy\color{fmA}base}: el factor que se repite};
  \node[fill=fmB!14,draw=fmB,rounded corners=2mm,font=\large,anchor=east] (le) at ([xshift=-1.6cm,yshift=0.8cm]p.west) {{\heavy\color{fmB}exponente}: cuántas veces};
  \draw[-{Latex[length=3mm]},fmA,line width=1.3pt] (lb.east) -- (pb);
  \draw[-{Latex[length=3mm]},fmB,line width=1.3pt] (le.east) to[out=0,in=150] (pe);
\end{tikzpicture}
\end{center}

\vfill
\begin{minipage}[t]{0.5\linewidth}
\begin{idea}[Regla de los signos]
Base \textbf{positiva} $\to$ siempre positivo.\\[2pt]
Base \textbf{negativa}: exponente \textbf{par} $\to$ \cK{positivo}; \textbf{impar} $\to$ \cR{negativo}.\\[4pt]
{\centering$(-2)^4=\cK{16}\qquad (-2)^3=\cR{-8}$\par}
\end{idea}
\end{minipage}\hfill
\begin{minipage}[t]{0.47\linewidth}
\begin{compte}[¡Cuidado con los paréntesis!]
\be{(-2)^4=16}\qquad \be{-2^4=-16}\\[3pt]
Sin paréntesis, el exponente afecta \textbf{solo al número}, no al signo.
\end{compte}
\end{minipage}

\vfill
\begin{tcbraster}[raster columns=3,raster equal height,raster column skip=8pt]
\begin{tcolorbox}[fmbase,colback=fmdefbg,colframe=fmdef,coltitle=white,title={Casos especiales},before upper=\centering]
$a^1=a$\par $a^2$: «al cuadrado»\par $a^3$: «al cubo»
\end{tcolorbox}
\begin{tcolorbox}[fmbase,colback=fmdefbg,colframe=fmdef,coltitle=white,title={De una fracción},before upper=\centering]
{\large$\left(\dfrac{a}{b}\right)^{n}=\dfrac{a^{n}}{b^{n}}$}\par\smallskip
$\left(\dfrac23\right)^{3}=\dfrac{8}{27}$
\end{tcolorbox}
\begin{tcolorbox}[fmbase,colback=fmdefbg,colframe=fmdef,coltitle=white,title={Potencias de 10},before upper=\centering]
$10^n$: un 1 y $n$ ceros\par\smallskip
$10^6=1\,000\,000$
\end{tcolorbox}
\end{tcbraster}

\vfill
\begin{exemple}[Ejemplos]
\begin{minipage}[t]{0.32\linewidth}
$(-5)^2=25$\\[3pt]
$-(-3)^3=-(-27)=27$
\end{minipage}\hfill
\begin{minipage}[t]{0.32\linewidth}
$(-1)^{7}=-1$ \ \nota{impar}\\[3pt]
$(-1)^{20}=1$ \ \nota{par}
\end{minipage}\hfill
\begin{minipage}[t]{0.32\linewidth}
$\left(-\dfrac25\right)^{3}=-\dfrac{8}{125}$\\[3pt]
$0^5=0\qquad 1^{37}=1$
\end{minipage}
\end{exemple}

\vfill
\begin{exemple}[Ejemplo: \ calcula]
\begin{minipage}[t]{0.19\linewidth}\centering $2^5=32$\end{minipage}\hfill
\begin{minipage}[t]{0.19\linewidth}\centering $(-3)^4=81$\end{minipage}\hfill
\begin{minipage}[t]{0.19\linewidth}\centering $-3^4=\cR{-81}$\end{minipage}\hfill
\begin{minipage}[t]{0.19\linewidth}\centering $\left(\dfrac12\right)^{4}=\dfrac{1}{16}$\end{minipage}\hfill
\begin{minipage}[t]{0.19\linewidth}\centering $10^3=1000$\end{minipage}
\end{exemple}

\vfill
\begin{truc}[¿Cómo se lee?]
$a^n$ se lee «$a$ elevado a $n$». \ $5^2$: «cinco al cuadrado». \ $2^3$: «dos al cubo». \ $7^4$: «siete elevado a cuatro».
\end{truc}

\vfill
% ============================================================ 1.10
\clearpage
\seccio{sec10}{Propiedades de las potencias}{Para multiplicar y dividir, jugamos con los exponentes.}

\providecommand{\prop}{}\renewcommand{\prop}[3]{\begin{tcolorbox}[fmbase,colback=fmdefbg,colframe=fmdef,coltitle=white,title={#1},before upper=\centering]
{\Large$#2$}\par\smallskip{\small$#3$}
\end{tcolorbox}}
\begin{tcbraster}[raster columns=3,raster equal height=rows,raster column skip=8pt]
\prop{\pas{1}\ Producto}{a^{m}\cdot a^{n}=a^{\cK{m+n}}}{2^{5}\cdot2^{3}=2^{8}=256}
\prop{\pas{2}\ Cociente}{a^{m}:a^{n}=a^{\cK{m-n}}}{7^{9}:7^{7}=7^{2}=49}
\prop{\pas{3}\ Potencia de potencia}{\left(a^{m}\right)^{n}=a^{\cK{m\cdot n}}}{\left(5^{2}\right)^{3}=5^{6}=15\,625}
\prop{\pas{4}\ Mismo exponente}{a^{n}\cdot b^{n}=(\cA{a\cdot b})^{n}}{2^{4}\cdot5^{4}=10^{4}=10\,000}
\prop{\pas{5}\ Mismo exponente}{a^{n}:b^{n}=(\cA{a:b})^{n}}{15^{3}:5^{3}=3^{3}=27}
\prop{Exponente cero}{a^{0}=1\ \ (a\neq0)}{a^{n}:a^{n}=a^{0}=1}
\end{tcbraster}

\vfill
\begin{compte}[¡Solo para productos y cocientes, no para sumas!]
\mal{2^{2}+2^{3}=2^{5}}\ \ {\small$4+8=12$, \ pero $2^5=32$}\qquad\quad
\mal{(2+3)^{2}=2^{2}+3^{2}}\ \ {\small$25$, \ pero $4+9=13$}
\end{compte}

\vfill
\subsection{Bases distintas: descompón en factores primos}
\begin{minipage}[c]{0.42\linewidth}\raggedright
\pas{1} \textbf{Descompón} cada base: \ $12=2^2\cdot3$, \ $8=2^3$\dots\\[6pt]
\pas{2} \textbf{Reparte} el exponente a cada factor: \ $12^{3}=2^{6}\cdot3^{3}$.\\[6pt]
\pas{3} \textbf{Agrupa} las potencias de la misma base y calcula.
\end{minipage}\hfill
\begin{minipage}[c]{0.55\linewidth}
\begin{exemple}[Ejemplo: \ $\dfrac{8^{2}\cdot4^{3}}{16^{2}}$]
$=\dfrac{\left(2^{3}\right)^{2}\cdot\left(2^{2}\right)^{3}}{\left(2^{4}\right)^{2}}=\dfrac{2^{6}\cdot2^{6}}{2^{8}}=\dfrac{2^{12}}{2^{8}}=2^{4}=\boxed{16}$
\end{exemple}
\end{minipage}

\vfill
\begin{exemple}[Ejemplo: \ $\dfrac{6^{3}\cdot9^{2}}{4\cdot27^{2}}$]
$=\dfrac{(\cA{2}\cdot\cB{3})^{3}\cdot\left(\cB{3}^{2}\right)^{2}}{\cA{2}^{2}\cdot\left(\cB{3}^{3}\right)^{2}}=\dfrac{\cA{2}^{3}\cdot\cB{3}^{3}\cdot\cB{3}^{4}}{\cA{2}^{2}\cdot\cB{3}^{6}}=\dfrac{\cA{2}^{3}\cdot\cB{3}^{7}}{\cA{2}^{2}\cdot\cB{3}^{6}}=\cA{2}^{3-2}\cdot\cB{3}^{7-6}=2\cdot3=\boxed{6}$
\end{exemple}

\vfill
\begin{exemple}[Ejemplo: \ $\dfrac{12^{2}\cdot18^{3}}{2^{5}\cdot27^{2}}$]
$=\dfrac{\left(2^{2}\cdot3\right)^{2}\cdot\left(2\cdot3^{2}\right)^{3}}{2^{5}\cdot\left(3^{3}\right)^{2}}=\dfrac{2^{4}\cdot3^{2}\cdot2^{3}\cdot3^{6}}{2^{5}\cdot3^{6}}=\dfrac{2^{7}\cdot3^{8}}{2^{5}\cdot3^{6}}=2^{2}\cdot3^{2}=\boxed{36}$
\end{exemple}

\vfill
% ============================================================ 1.11
\clearpage
\seccio{sec11}{Potencias de exponente cero y negativo}{Un exponente negativo «gira» la base.}

\vfill
\begin{minipage}[c]{0.5\linewidth}
¿De dónde salen? \ $\dfrac{2^3}{2^3}=1$ y también $2^{3-3}=2^0$: por eso $\boxed{a^0=1}$.\\[4pt] Por la propiedad del cociente, \ $2^{3}:2^{5}=2^{-2}$.\\ Y si simplificamos:\par\smallskip
{\centering$\dfrac{2^{3}}{2^{5}}=\dfrac{\cancel{2}\cdot\cancel{2}\cdot\cancel{2}}{\cancel{2}\cdot\cancel{2}\cdot\cancel{2}\cdot2\cdot2}=\dfrac{1}{2^{2}}$\par}\smallskip
Por tanto, \ $2^{-2}=\dfrac{1}{2^{2}}$.
\end{minipage}\hfill
\begin{minipage}[c]{0.46\linewidth}
\begin{tcolorbox}[enhanced,colback=fmyellow!18,colframe=fmgold,arc=4mm,boxrule=1.6pt,before upper=\centering,top=3mm,bottom=3mm]
{\Large$a^0=1$}\par\medskip{\Large$a^{-n}=\dfrac{1}{a^{n}}$}\par\medskip
{\Large$\left(\dfrac{a}{b}\right)^{-n}=\left(\dfrac{b}{a}\right)^{n}$}
\end{tcolorbox}
\end{minipage}

\vfill
\begin{compte}[El resultado NO se vuelve negativo]
El exponente negativo \textbf{invierte} la base, no cambia el signo: \qquad \be{2^{-3}=\dfrac18}\qquad \mal{2^{-3}=-8}
\end{compte}

\vfill
\begin{exemple}[Ejemplos]
\begin{minipage}[t]{0.32\linewidth}
$2^{-3}=\dfrac{1}{2^{3}}=\dfrac18$\\[5pt]
$(-3)^{-2}=\dfrac{1}{(-3)^{2}}=\dfrac19$
\end{minipage}\hfill
\begin{minipage}[t]{0.32\linewidth}
$\left(\dfrac25\right)^{-2}=\left(\dfrac52\right)^{2}=\dfrac{25}{4}$\\[5pt]
$\left(-\dfrac12\right)^{-3}=(-2)^{3}=-8$
\end{minipage}\hfill
\begin{minipage}[t]{0.32\linewidth}
$10^{-4}=\dfrac{1}{10\,000}=0{,}0001$\\[5pt]
Y al revés: \ $\dfrac18=2^{-3}$
\end{minipage}
\end{exemple}

\vfill
\begin{idea}[Las propiedades también valen con exponentes negativos]
Cuidado con los signos al sumar y restar exponentes:\\[4pt]
{\centering$3^{-2}\cdot3^{5}=3^{-2+5}=3^{3}\qquad 5^{-3}:5^{-7}=5^{-3-(-7)}=5^{4}\qquad \left(2^{-3}\right)^{-2}=2^{6}\qquad 6^{-2}:3^{-2}=2^{-2}=\dfrac14$\par}
\end{idea}

\vfill
\begin{minipage}[t]{0.49\linewidth}
\begin{exemple}[Ejemplo: \ $\dfrac{4^{-2}\cdot8^{3}}{16^{-1}\cdot2^{5}}$]
$=\dfrac{2^{-4}\cdot2^{9}}{2^{-4}\cdot2^{5}}=\dfrac{2^{5}}{2^{1}}=2^{4}=\boxed{16}$
\end{exemple}
\end{minipage}\hfill
\begin{minipage}[t]{0.48\linewidth}
\begin{exemple}[Ejemplo: \ $\left(\dfrac23\right)^{-2}\cdot\left(\dfrac49\right)^{3}$]
$=\dfrac{3^{2}}{2^{2}}\cdot\dfrac{2^{6}}{3^{6}}=\dfrac{2^{4}}{3^{4}}=\boxed{\dfrac{16}{81}}$
\end{exemple}
\end{minipage}

\vfill
\begin{exemple}[Ejemplo: \ $\dfrac{6^{-2}\cdot9^{3}}{4^{-1}\cdot3^{5}}$]
$=\dfrac{(2\cdot3)^{-2}\cdot\left(3^{2}\right)^{3}}{\left(2^{2}\right)^{-1}\cdot3^{5}}=\dfrac{2^{-2}\cdot3^{-2}\cdot3^{6}}{2^{-2}\cdot3^{5}}=\dfrac{2^{-2}\cdot3^{4}}{2^{-2}\cdot3^{5}}=2^{0}\cdot3^{-1}=\boxed{\dfrac13}$
\end{exemple}

\vfill
% ============================================================ resum
\resumtitol
\begin{tcbraster}[raster columns=2,raster equal height=rows,raster column skip=14pt,raster row skip=26pt,raster before skip=16pt]
\rtile{1.1}{Naturales}{$\mathbb N=\{0,1,2,\dots\}$. \ $D=d\cdot q+r$, \ $r<d$.}
\rtile{1.2}{Jerarquía}{Paréntesis, potencias y raíces, $\cdot$ y $:$, $+$ y $-$.}
\rtile{1.3}{Enteros}{$|a|$: distancia al 0. \ Más a la derecha, mayor.}
\rtile{1.4}{Operaciones con enteros}{Restar es sumar el opuesto. Signos iguales, $+$; distintos, $-$.}
\rtile{1.5}{Fracciones}{Equivalentes: $a\cdot d=b\cdot c$. Simplificamos por el m.c.d.}
\rtile{1.6}{Operaciones}{$\pm$: común denominador. \ $:$ multiplicar por la inversa.}
\rtile{1.7}{Combinadas}{La misma jerarquía. Simplifica cada paso.}
\rtile{1.8}{Problemas}{Fracción de una cantidad y de la parte que queda.}
\rtile{1.9}{Base negativa}{Exponente par, $+$; impar, $-$. \ $-3^2\neq(-3)^2$.}
\rtile{1.10}{Propiedades}{$a^m\cdot a^n=a^{m+n}$, \ $a^m:a^n=a^{m-n}$, \ $(a^m)^n=a^{mn}$.}
\rtile{1.11}{Exponente negativo}{$a^0=1$, \ $a^{-n}=\frac1{a^n}$, \ $\left(\frac ab\right)^{-n}=\left(\frac ba\right)^n$.}
\end{tcbraster}
\end{document}
